\chapter*{Preface: The Road That Listens}

A story is not a reckoning. It balances nothing. But it warms the hands.

\smallskip

You travel. You sit by fires. A woman speaks. You write. The ash on your cuff becomes ink. That is all. That is enough.

\smallskip

My people are the Tulkani. We build no temples. Our wagons carry the hearth. Kuva rides the storm—some say. I have lit a candle for him when the road was ice and the horses would not move. Kuva counts the stars. I count the thresholds.

\smallskip

Ikasha does not command. She sleeps. Between her breaths there is room for secrets. I have borrowed that space. I will return it when I am done. She does not ask for interest.

\smallskip

These tales were told by midwives, bell-wardens, hedge-mothers, cord-weavers. They did not call themselves witches. The women who speak to the dead never do. The dead have enough names already.

\smallskip

I wrote as the fire guttered. Where a teller coughed, I left a space. Where a bell rang, I marked the hour. Where the ninth word was missing—I did not add it. The ninth word belongs to the threshold. The threshold was not speaking that night.

\smallskip

Some tales are comfort. Some are warnings. A few are keys. The lock is not in this book. The lock is in the reader's chest. Turn the key gently. It may be older than you think.

\smallskip

Read them. Or do not. The road turns regardless. The hearth waits. The stories will be here when you come back. They have nowhere else to go.

\smallskip
\noindent --- \textit{Dusana of the Raven Road}

\noindent \textit{Under the lamp, bread and salt. The ink is ash and rain. The rain has stopped. The ash is still warm.}